\scanpage{646}
\mspage{3\deleted{\unsure{5}}\unsure{4|5}}

Alors, par le th. de factorisation 1.1, on a

\[
 C\mathrel{\leftarrow\!\!|\!\!\rightarrow}TF\overset{\mathrm{def}}{=}C_*,
\]

et l'axiome de factorisation $\mathrm{Cof}_3$ implique
(cf 1 \unsure{M|11}) $TF\subset W^{\unsure{\natural|\sharp}}$. Plus généralement\unsure{.}

\uline{Prop. 3.1} Supposons \inserted{que $M$ soit \unsure{spécial|spéciale} (stable par petites $\varinjlim$, et éléments accessibles),} que \deleted{\unsure{f}} $(W,C)$ satisfasse
à $\mathrm{Cof}_0$, $\mathrm{Cof}_1$, et que $C$ soit stable par
composition transfinie (ce qui, joint à $\mathrm{Cof}_1$,
\unsure{s'énonce|se résume} en disant que $C$ est $Q$-saturé
transfini à gauche), enfin que $C$
admette une petite partie $C_0$, $Q$-génératrice
à gauche, de sorte que $C\mathrel{\leftarrow\!\!|\!\!\rightarrow}TF\overset{\mathrm{def}}{=}C_*$.

\deleted{Supposons qu'on ait $W=W^{\natural}$, i.e. $W$ stable par
facteurs directs}. Alors, \inserted{pour que} l'axiome de
factorisation $\mathrm{Cof}_3$ \deleted{équivaut à $TF\subset W$}\inserted{soit satisfait, il suffit qu'on ait}

\[
 TF\subset W,\tag*{(3.1.1)}
\]

et il faut qu'on ait $TF\subset W^{\natural}$ ; donc si $W=W^{\natural}$,
l'axiome $\mathrm{Cof}_3$ équivaut à (3.1.1).

La structure de cofibrations est dite \inserted{\uline{semi}-}\uline{close}
si \inserted{$C$ est $Q$-clos à g., i.e. si} on a $C=\unsure{C^{\sim}|\widetilde C}$, (ce qui implique
que $C$ est stable par composition transfinie),
\uline{close} si de plus $TC=(TC)^{\sim}$. (ce qui
implique l'axiome transfini $\mathrm{Cof}_4$).
% archive sheet number: 645
